
\begin{proposition}[Refuted]\label{cdp:st:D.1.5}

\textbf{Explicit second examples for the erroneous steps.}
Let $E_0=F_0=E'=\C/(\Z+i\Z)$ and $a=1/2\in E_0$.
The involution
\[
(e,z)\longmapsto(e+a,-z)
\]
acts freely on $E_0\times F_0$. Put
$S=(E_0\times F_0)/\langle(a,-1)\rangle$,
$Y=S\times E'$, and $C=E_0/\langle a\rangle$.
The induced map $Y\to C$ is a smooth torus fibration with fibre
$F_0\times E'$ and monodromy
$\operatorname{diag}(-I_2,I_2)$ around the nontrivial deck character.
Its integral homology coinvariants are
$\Z^2\oplus(\Z/2\Z)^2$. After removing $s\geq1$ fibres,
Proposition \ref{cdp:st:D.1.1} gives
\[
H_1(U;\Z)=\Z^{s+3}\oplus(\Z/2\Z)^2,
\]
whereas the uncorrected direct-sum argument has free rank $s+5$.
This is an explicit compact complex realization of the faulty
group-theoretic step, independent of the retained six-sphere.

Let $\eta$ be the nontrivial order-two degree-zero line bundle on
$C$ associated to the double cover $E_0\to C$. Finite-cover descent
of fibre cohomology gives
\[
R^1f_*\OO_Y=\eta\oplus\OO_C,\qquad R^2f_*\OO_Y=\eta.
\]
The invariant part of the exterior algebra on the three
antiholomorphic one-form directions, with signs $(1,-1,1)$,
gives $h^q(Y,\OO_Y)=(1,2,1,0)$. The restriction from
$H^1(Y,\OO_Y)$ to the two-dimensional $H^1(F_0\times E',\OO)$
has image precisely the $E'$ direction and has rank one.
All these homomorphisms are induced by the displayed finite
\'etale quotient and projection. This example has $b_1(Y)=4$;
it is a counterexample to the unchecked restriction argument, not
a second example satisfying the integral CDP20 hypotheses.

There is also an explicit negative summand over $\PP^1$. Let
$h:R\to\PP^1$ be the relatively minimal elliptic surface with
Weierstrass equation
\[
y^2=x^3+(t^4+1)x+t^6.
\]
Its fundamental line bundle is $\OO(1)$ and its discriminant has
twelve simple zeroes. Relative duality gives
$R^1h_*\OO_R=\OO(-1)$. For $Y'=R\times E'\to\PP^1$,
the product K\"unneth maps give
\[
R^1f_*\OO_{Y'}=\OO\oplus\OO(-1),\qquad
R^2f_*\OO_{Y'}=\OO(-1),
\]
and the restriction in degree one again has rank one. This
projective example has $b_1=2$, $b_2=11$ and $a=3$, so it does not
replace the missing second counterexample under all the standing
CDP20 hypotheses.
\end{proposition}

The bielliptic-product rank defect is already described in \citet[Section~10.5, p.~91]{AC2026}. The full integral torsion calculation, the coherent direct images and the separate rational-elliptic product below retain their additional precise content.



